\section{Vyhodnotenie podľa princípov použiteľnosti}

Používateľské rozhranie bolo analyzované aj z pohľadu základných princípov použiteľnosti, konkrétne podľa Nielsenových heuristík, ktoré predstavujú jeden z najčastejšie používaných prístupov pri hodnotení používateľských rozhraní \cite{nielsen}. 
V kontexte informačných systémov zohráva použiteľnosť významnú úlohu pri efektívnom využívaní systému a ovplyvňuje produktivitu používateľov \cite{laudon}.

\subsubsection*{1. Viditeľnosť stavu systému}

Systém by mal používateľa priebežne informovať o tom, čo sa deje.

\textbf{Hodnotenie:} Čiastočne splnené.  
Pri niektorých operáciách (napr. vytvorenie záznamu alebo spisu) nie je vždy zrejmé, či operácia prebehla úspešne. Chýba výraznejšia spätná väzba (napr. notifikácia alebo zvýraznenie zmeny).

\subsubsection*{2. Zhoda medzi systémom a reálnym svetom}

Systém by mal používať pojmy a koncepty zrozumiteľné používateľovi.

\textbf{Hodnotenie:} Čiastočne splnené.  
Niektoré pojmy ako „vecná skupina“, „evidencia“ alebo „proces“ nemusia byť pre nového používateľa jednoznačné bez predchádzajúcej znalosti domény.

\subsubsection*{3. Kontrola a sloboda používateľa}

Používateľ by mal mať možnosť jednoducho vrátiť alebo opraviť svoje kroky.

\textbf{Hodnotenie:} Slabo splnené.  
Rozhranie neposkytuje explicitné možnosti návratu späť alebo zrušenia operácií, čo môže viesť k frustrácii pri chybných krokoch.

\subsubsection*{4. Konzistentnosť a štandardy}

Používateľ by nemal byť nútený premýšľať, či rôzne prvky znamenajú to isté.

\textbf{Hodnotenie:} Splnené.  
Tabuľky, formuláre a navigačné prvky majú konzistentný vzhľad a správanie naprieč celou aplikáciou.

\subsubsection*{5. Prevencia chýb}

Systém by mal predchádzať vzniku chýb.

\textbf{Hodnotenie:} Čiastočne splnené.  
Formuláre obsahujú validačné mechanizmy, avšak chýbajú vysvetlenia alebo odporúčania, ktoré by používateľa naviedli k správnemu vyplneniu.

\subsubsection*{6. Rozpoznanie namiesto zapamätania}

Používateľ by nemal byť nútený pamätať si informácie medzi krokmi.

\textbf{Hodnotenie:} Slabo splnené.  
Používateľ si musí pamätať, kde sa nachádzajú jednotlivé funkcie (napr. spustenie procesu), čo znižuje intuitívnosť systému.

\subsubsection*{7. Flexibilita a efektívnosť používania}

Systém by mal umožniť efektívnu prácu skúseným používateľom.

\textbf{Hodnotenie:} Čiastočne splnené.  
Pre skúseného používateľa je práca relatívne efektívna, avšak chýbajú skratky alebo zrýchlené postupy.

\subsubsection*{8. Estetický a minimalistický dizajn}

Rozhranie by nemalo obsahovať zbytočné informácie.

\textbf{Hodnotenie:} Čiastočne splnené.  
Niektoré obrazovky (najmä formuláre) obsahujú veľké množstvo polí, čo môže pôsobiť neprehľadne.

\subsubsection*{9. Pomoc pri rozpoznaní, diagnostike a oprave chýb}

Chybové hlášky by mali byť zrozumiteľné.

\textbf{Hodnotenie:} Nedostatočne splnené.  
Systém neposkytuje dostatočne detailné informácie o chybách ani návrhy na ich riešenie.

\subsubsection*{10. Pomoc a dokumentácia}

Systém by mal poskytovať pomoc používateľovi.

\textbf{Hodnotenie:} Nesplnené.  
V aplikácii chýba integrovaná pomoc, onboarding alebo dokumentácia priamo v rozhraní.

\subsubsection*{Zhrnutie hodnotenia}

Na základe vykonanej analýzy možno konštatovať, že používateľské rozhranie je funkčné a konzistentné, avšak vykazuje nedostatky najmä v oblasti používateľskej podpory, navigácie a spätnej väzby.

Najvýznamnejšie nedostatky boli identifikované v oblastiach:
\begin{itemize}
    \item viditeľnosť stavu systému,
    \item používateľské vedenie (guidance),
    \item pomoc a dokumentácia.
\end{itemize}

Tieto nedostatky môžu mať negatívny vplyv najmä na menej skúsených používateľov a znižovať celkovú použiteľnosť systému.
